\documentclass[11pt]{amsart}
\usepackage[a4paper,margin=1in]{geometry}
\raggedbottom
\usepackage[bottom]{footmisc}
\usepackage{amsmath,amssymb,amsthm}
\usepackage{booktabs}
\usepackage{needspace}
\usepackage[hidelinks]{hyperref}
\hypersetup{pdftitle={Explicit initial Theta-data and the two readings of Step (v) of IUT IV, Theorem 1.10: a catalogue},
  pdfauthor={Toshihisa Urahigashi}}

\newtheorem{theorem}{Theorem}
\newtheorem{proposition}[theorem]{Proposition}
\newtheorem{lemma}[theorem]{Lemma}
\theoremstyle{remark}
\newtheorem{remark}[theorem]{Remark}

\newcommand{\Q}{\mathbb Q}
\newcommand{\Z}{\mathbb Z}
\newcommand{\F}{\mathbb F}
\newcommand{\Qp}{\mathbb Q_p}
\newcommand{\cO}{\mathcal O}
\newcommand{\cA}{\mathcal A}
\newcommand{\dd}{\mathfrak d}
\newcommand{\qq}{\underline{\underline{q}}{}}
\newcommand{\Fmod}{F_{\mathrm{mod}}}
\newcommand{\VV}{\underline{\mathbb V}}
\newcommand{\Vbadmod}{\mathbb V^{\mathrm{bad}}_{\mathrm{mod}}}
\newcommand{\Vbad}{\underline{\mathbb V}^{\mathrm{bad}}}
\newcommand{\mulog}{\mu^{\log}}
\newcommand{\RIt}{(R_I)^{\sim}}
\newcommand{\pr}{\operatorname{pr}}
\newcommand{\Nm}{\operatorname{N}}
\newcommand{\Sj}{\mathbb S^{\pm}_{j+1}}
\newcommand{\bv}{\mathbf v}
\newcommand{\bu}{\mathbf u}

\begin{document}

\title[Explicit initial $\Theta$-data and Step~(v): a catalogue]{Explicit initial $\Theta$-data and the two readings of Step~(v) of IUT~IV, Theorem~1.10: a catalogue}
\author{Toshihisa Urahigashi}
\email{123026.urahigashi.toshihisa@gmail.com}
\date{Version 1.2.0, 8 November 2026}

\begin{abstract}
We collect the explicit initial $\Theta$-data, in the sense of Definition~3.1 of the first paper of Mochizuki's
series on inter-universal Teichm\"uller theory, that the author has constructed over $\Q(\sqrt5)$ and $\Q(i)$:
three published earlier, three given here, and one, for fourteen values of $\ell$, that is verified in a separate
note. Each has a split rational
prime above which the two places of the field of moduli are either one of good and one of multiplicative reduction,
or both of multiplicative reduction with different orders of the $q$-parameter. For each datum we tabulate the local
invariants and compare, in Step~(v) of the proof of Theorem~1.10 in the fourth paper, the displayed comparison module
with the hull component of the union of possible images. Under the reading in which the theta value stays at the
label $j$, as in Step~(v) and in what the public Lean formalization lana-agents/iut calls the reading of IUT~IV,
every comparison is a containment. Under the reading in which label permutations fixing the label $0$ act before
the hull is taken, which is contained in what the same formalization calls the literal reading (that reading also
admits permutations moving $0$), and under stated membership hypotheses, every datum gives a
non-containment for a single component; at primes with a good and a bad place the excess grows like $j^2$, and for
three data the weighted average over the components also exceeds the bound displayed at the prime. No conclusion is
drawn about Corollary~3.12, the final inequality of Theorem~1.10, or the abc conjecture.
\end{abstract}

\maketitle

\section{Introduction}

Explicit initial $\Theta$-data in the sense of \cite[Def.~3.1]{IUT1} are needed whenever the estimates of
\cite{IUT4} are to be evaluated on actual numbers. Dupuy and Hilado computed one for the curve 11a1 over $\Q$
\cite[\S2.10]{DH1}, and Zhou constructed families from Frey--Hellegouarch curves \cite{Zhou1,Zhou2}; in these
constructions the field of moduli is $\Q$, so a rational prime carries a single place of it. The author has published
three data over quadratic fields \cite{First,Second,Note}. This note collects them with three further data and with
the datum of \cite{Fourteen}, given for fourteen values of $\ell$, all over $\Q(\sqrt5)$ or $\Q(i)$, and records for
each the local invariants above one split prime and the comparisons in
Step~(v) of \cite[Thm.~1.10, proof]{IUT4}.

The comparisons depend on how the label permutations contained in the indeterminacy (Ind1) enter the union of
possible images of \cite[Cor.~3.12]{IUT3}. That they permute tensor factors, and that a simplified computation of the
hull may omit this, is noted in \cite[\S6.2, footnote~7]{DH}; procession automorphisms are described as permutations
of tensor factors in \cite[\S4.6]{DH1}. The public Lean repository \cite{LANAiut} formalizes two readings of
Corollary~3.12, one with the theta value at the label $j$ (its reading of IUT~IV) and one with the theta value at any
label (its literal reading), shows that the first implies the second, and records in prose that the two give the same
regions when the orders of the $q$-parameters agree at all places above each prime, that two bad places with different
orders already break this, and that at primes with a good and a bad place, which arise from the Belyi reduction, the
literal reading is governed by the minimum rather than the average. The data below are explicit instances of these
cases.

What is new here is the data themselves and the numbers: initial $\Theta$-data whose field of moduli is not $\Q$
(we have not found such data in the literature; an explicit example of this kind was asked for in \cite{MO}), data
realizing both kinds of prime just described, and exact values of the comparisons of Step~(v) on them. The
distinctions between the readings, and the mechanism by which the readings differ, are not new; see \S\ref{s:prior}.

We do not claim that Corollary~3.12, the final inequality of Theorem~1.10 or the abc conjecture is false, that either
reading is the one the source intends, or that the hypotheses of \S\ref{s:readings} identify the source's union of
possible images.

\section{Step (v) and the readings}\label{s:readings}

We use the notation of \cite[\S1]{Note}. For a rational prime $p$ and a place $v$ of $\Fmod$ above $p$ let $k_v$ be
the completion of $K$ at the selected place above $v$, $e_v$ its ramification index over $\Qp$, $\dd_v=(e_v-1)/e_v$
the order of its different when $p\nmid e_v$, and $\rho(v)=\operatorname{ord}_p(\qq_v)$ if $v\in\Vbadmod$,
$\rho(v)=0$ otherwise. For $j\in\{1,\dots,\ell^*\}$ and a tuple $\bv=(v_0,\dots,v_j)$ of places above $p$, Step~(v)
compares the component at $\bv$ of the union $U_j$ of possible images with a module $M_{\bv}$; when every $e_{v_i}\le
p-2$,
\begin{equation}\label{eq:bound}
\mulog(M_{\bv})\le\Bigl(-j^2\rho(v_j)+\dd_I+1\Bigr)\log p,\qquad \dd_I=\textstyle\sum_{i=0}^{j}\dd_{v_i}
\end{equation}
\cite[Prop.~1.4(iii), p.~13]{IUT4}, \cite[(1)]{Note}. All values below are in units of $\log p$.

We consider the following readings of \cite[Cor.~3.12]{IUT3}; none is proved here.
\begin{itemize}
\item[(V)] The theta value stays at the label $j$, as in Step~(v) and in the reading of IUT~IV of \cite{LANAiut}.
(The file \nolinkurl{Iut/Concrete/Ind1Invariance.lean} of \cite{LANAiut} also records that taking the hull for each
image under (Ind1) separately gives the value of IUT~IV.)
\item[(L)] Label permutations fixing $0$ act on possible-image regions before the hull is taken; this is (R) of
\cite[\S4]{Note}. (L$^0$) also admits permutations moving $0$, as in \cite[(SA2)]{FirstFull}. That each stage of a
procession of copies of one prime-strip carries the full symmetric group, so that the fixing of $\mathcal D_0$ does not
by itself fix the label $0$ inside the larger capsules, is shown in \cite[\S6]{Fourth} from \cite[\S0, Def.~4.10,
Prop.~6.9(i)]{IUT1}.
\item[(M$_{\mathrm b}$)] At a bad place the region of the $\Theta$-pilot object is the theta value times a lattice
containing $1$, and its Kummer image is a possible image; this is (M) of \cite[\S4]{Note}, which gives
\cite[Lemma~3]{Note}.
\item[(M$_{\mathrm g}$)] At a good place the region of the $\Theta$-pilot object at a tuple whose factor of label $j$
lies over it contains $1\otimes\cdots\otimes1$, the theta value being $1$ there, and its Kummer image is a possible
image; a source basis for this is given in \cite[\S7, (SA1)]{FirstFull}. The proof of
\cite[Lemma~3]{Note} then gives, at such a tuple $\bu$, an element $X\in U_j$ with $\pr_{\bu}(X)=\chi(1)$,
$\chi=\otimes_i\chi_i$, $\chi_i\in\cA(k_{u_i})$.
\end{itemize}
Under (V), for every datum and every tuple, the elements of \cite[Lemma~3]{Note} lie in $M_{\bv}$: this is
\cite[Prop.~6]{Note} with $b=v_j$, and at a good $v_j$ the same argument with $\lambda=0$ applies to $\chi(1)$.

Under (L), elements are moved by \cite[Lemma~4]{Note}. We use three kinds, with lower bounds for the hull component
obtained from \cite[(1.1)--(1.3)]{Note} and the argument of \cite[Thm.~5]{Note} (if $U_j$ is not relatively compact
its hull has log-volume $+\infty$ and every finite bound fails).
\begin{itemize}
\item[(BB)] Two bad places $b_s,b_l$ above $p$ with $\rho(b_s)<\rho(b_l)$, $j\ge2$, assuming (M$_{\mathrm b}$) and
(L): the theta value of $b_s$ is moved by a transposition $(t\ j)$, $1\le t<j$, to a tuple with $v_t=b_s$ and
$v_j=b_l$; the hull component is $\ge-\lceil j^2\rho(b_s)\rceil$.
\item[(GB)] A good place $g$ and a bad place $b$ above $p$, assuming (M$_{\mathrm g}$) and (L) if $j\ge2$, (L$^0$)
if $j=1$: the element $\chi(1)$ at a tuple with label $j$ over $g$ is moved to a tuple with label $j$ over $b$; the
hull component is $\ge0$.
\item[(GB$_L$)] The variant of \cite[Thm.~8.1]{FirstFull}, under its hypothesis (SA): the moved region contains the
log-unit lattice, and the hull component is $\ge-\sum_i e_{v_i}^{-1}$. The same comparison is stated in
\cite[Thm.~5]{First} for $j=1$ and, for every $j$, under the hypothesis that the hull contains the log-unit lattice,
in \cite[Prop.~5]{Fourth}.
\end{itemize}

\begin{proposition}\label{p:closed}
Let $j\ge1$ and let $\bv$ range over the tuples of places above $p$. Among the elements just described, the largest
excess of the hull component over the bound \eqref{eq:bound} is
\begin{align*}
\text{(BB)}\quad & j^2\rho(b_l)-\lceil j^2\rho(b_s)\rceil-1-\dd_{b_s}-\dd_{b_l}-(j-1)\,\dd_{\min},
 &&\bv=(b_*,b_s,b_*,\dots,b_*,b_l),\\
\text{(GB)}\quad & j^2\rho(b)-\dd_b-1, &&\bv=(g,\dots,g,b),\\
\text{(GB$_L$)}\quad & j^2\rho(b)-(j+2), &&\text{any }\bv\text{ with }v_j=b,
\end{align*}
where $\dd_{\min}$ is the least $\dd_v$ over the places above $p$ and $b_*$ a place attaining it. A positive value is a
non-containment of the component at $\bv$ in $M_{\bv}$.
\end{proposition}

\begin{proof}
The bound \eqref{eq:bound} at $\bv$ is $-j^2\rho(v_j)+\sum_i\dd_{v_i}+1$. In (BB) and (GB) the lower bound does not
depend on the factors other than the moved one and the one of label $j$, so their $\dd$ is minimized; in (GB) the good
place has $\dd=0$. In (GB$_L$), $\sum_i(e_{v_i}^{-1}+\dd_{v_i})=j+1$ for tame $e_{v_i}$. The values were also checked
against an enumeration of all tuples for $\ell=7$ and $j\le3$ (\S\ref{s:repro}).
\end{proof}

\section{The data}\label{s:data}

All data have $E_0\colon y^2=x(x-1)(x-a)$ over $F_0=\Fmod$, $F=F_0(\sqrt{-1},E_0[15])$, $K=F(E_0[\ell])$, and
$\Vbadmod$ the places of multiplicative reduction outside $2\ell$, as in \cite[\S2]{Note}; $\pi$ generates a prime of
$F_0$ of degree one over $p=\Nm(\pi)$, and $\bar\pi$ is its conjugate. Table~\ref{tab:data} lists them.

\begin{table}[ht]
\centering\footnotesize\setlength{\tabcolsep}{3pt}
\begin{tabular}{@{}lllllll@{}}
\toprule
datum & $\pi$ & $p$ & $a$ & $\ell$ & $(e,\rho)$ above $p$ & proof\\
\midrule
$D_{211}$ & $16+3\sqrt5$ & $211$ & $\pi^{29}$ & $7$ & $g\,(1,0)$, $b\,(105,\tfrac{29}7)$ & \cite[Thm.~1]{First}\\
$D_{109}$ & $10+3i$ & $109$ & $\pi^{29}$ & $7$ & $g\,(1,0)$, $b\,(105,\tfrac{29}7)$ & \cite{Second}\\
$D_2$ & $14164+5825\sqrt5$ & $30965771$ & $\pi\bar\pi^{6}$ & $7$ & $b_1(105,\tfrac17)$, $b_2(35,\tfrac67)$ & \cite[Prop.~2]{Note}\\
$D_3$ & as $D_2$ & $30965771$ & $\pi^{6}\bar\pi^{-30}$ & $7$ & $b_1(35,\tfrac67)$, $b_2(7,\tfrac{30}7)$ & Prop.~\ref{p:D3}\\
$D_{18}$ & as $D_2$ & $30965771$ & $\pi^{18}$ & $7$ & $g\,(1,0)$, $b\,(35,\tfrac{18}7)$ & Prop.~\ref{p:D18}\\
$D_{15}$ & $-5932-919\sqrt5$ & $30965819$ & $\pi^{15}$ & $7$ & $g\,(1,0)$, $b\,(7,\tfrac{15}7)$ & Prop.~\ref{p:D15}\\
$D_{2357}^{(\ell)}$ & $41+26i$ & $2357$ & $\pi^{472}$ & $\ell\in\Lambda$ & $g\,(1,0)$, $b\,(15\ell,\tfrac{472}\ell)$ & \cite{Fourteen}\\
\bottomrule
\end{tabular}
\caption{The data; $\pi$ lies in $\Z[\sqrt5]$ or $\Z[i]$, and $F_0=\Q(\pi)$. $\Lambda=\{89,97,101,103,107,109,113,127,131,137,139,149,151,157\}$. In every case $e\le p-2$ at
both places, so \eqref{eq:bound} applies. The curve of $D_{2357}^{(\ell)}$ is that of \cite[\S13]{Second}; the
verification of \cite[Def.~3.1]{IUT1} for every $\ell\in\Lambda$ is in \cite{Fourteen}.}\label{tab:data}
\end{table}

\begin{proposition}\label{p:new}
The data $D_3$, $D_{18}$ and $D_{15}$ extend to initial $\Theta$-data in the sense of \cite[Def.~3.1]{IUT1} satisfying
the hypotheses of \cite[Thm.~1.10]{IUT4}, with $\Fmod=F_0$ and the local data of Table~\ref{tab:data}; for each,
$p>e^*_{\mathrm{mod}}\cdot\ell=7741440$.
\end{proposition}

The proofs are in Appendix~\ref{s:proofs}; each is the proof of \cite[Prop.~2]{Note} with the changes stated there.

\section{The comparisons}\label{s:comparisons}

\begin{theorem}\label{t:component}
For each positive entry of Table~\ref{tab:component}, under the readings stated for its column, the component of $U_j$
at the tuple of Proposition~\ref{p:closed} is not contained in $M_{\bv}$, and the log-volume of its hull exceeds the
bound \eqref{eq:bound} by at least the value shown. Under (V) the elements considered lie in $M_{\bv}$ for every
datum and every tuple.
\end{theorem}

\begin{proof}
Proposition~\ref{p:closed} with the values of Table~\ref{tab:data}, and the remark on (V) in \S\ref{s:readings}.
\end{proof}

\begin{table}[ht]
\centering\small
\begin{tabular}{@{}lllll@{}}
\toprule
datum & kind & $j=1$ & $j=2$ & $j=3$\\
\midrule
$D_2$ & (BB) & --- & $-\tfrac{158}{105}$ & $\tfrac{17}{21}$ \ (\cite[Thm.~5]{Note}: $\tfrac{27}{35}$)\\
$D_3$ & (BB) & --- & $\tfrac{331}{35}$ & $\tfrac{911}{35}$\\
$D_{211}$, $D_{109}$ & (GB) / (GB$_L$) & $\tfrac{226}{105}$ / $\tfrac87$ & $\tfrac{1531}{105}$ / $\tfrac{88}7$ & $\tfrac{3706}{105}$ / $\tfrac{226}7$\\
$D_{18}$ & (GB) / (GB$_L$) & $\tfrac35$ / $-\tfrac37$ & $\tfrac{291}{35}$ / $\tfrac{44}7$ & $\tfrac{741}{35}$ / $\tfrac{127}7$\\
$D_{15}$ & (GB) / (GB$_L$) & $\tfrac27$ / $-\tfrac67$ & $\tfrac{47}7$ / $\tfrac{32}7$ & $\tfrac{122}7$ / $\tfrac{100}7$\\
\bottomrule
\end{tabular}
\caption{Largest excess of the hull component over \eqref{eq:bound}, in units of $\log p$. (BB) assumes (M$_{\mathrm b}$)
and (L); (GB) assumes (M$_{\mathrm g}$), and (L) for $j\ge2$, (L$^0$) for $j=1$; (GB$_L$) assumes (SA) of
\cite{FirstFull}. For $D_2$, $j=3$, the tuple of Proposition~\ref{p:closed} is $(b_2,b_1,b_2,b_2)$; the tuple
$(b_1,b_1,b_1,b_2)$ of \cite[Thm.~5]{Note} gives $\tfrac{27}{35}$.}\label{tab:component}
\end{table}

For $D_{2357}^{(\ell)}$ the values are $\tfrac{472j^2}\ell-2+\tfrac1{15\ell}$ (GB) and $\tfrac{472j^2}\ell-(j+2)$
(GB$_L$), positive for every $j\ge1$ and $\ell\in\Lambda$; at $\ell=157$ they are $\tfrac{2371}{2355}$ and
$\tfrac1{157}$ for $j=1$, and about $18288.8$ and $18210.8$ for $j=\ell^*=78$. At primes with a good and a bad place the
excess grows like $j^2\rho(b)$, the order of the term $-j^2\rho(v_j)$ of \eqref{eq:bound} itself; at the bad-bad
primes of $D_2$ and $D_3$ it is of the order $j^2(\rho(b_l)-\rho(b_s))$.

\begin{theorem}\label{t:average}
For $D_2$, $D_3$, $D_{18}$ and $D_{15}$ the two places above $p$ have degree one, so the tuples have equal weights, and
$p>e^*_{\mathrm{mod}}\ell$, so the terms with $\log(\mathfrak s^{\le}_p)$ vanish \cite[pp.~23, 25]{IUT4}. Under
(M$_{\mathrm b}$), (M$_{\mathrm g}$) and (L), the weighted average over the tuples of the hull components at stage $j$ is
at least $L_j=-\mathbb E_{\bv}\bigl[\min_{1\le t\le j}c(v_t)\bigr]$, where $c(b)=\lceil j^2\rho(b)\rceil$ for a bad
place and $c(g)=0$ for a good one. The weighted average that Step~(v) displays at stage $j$ is
$B_j=(j+1)\bar\dd-j^2\bar\rho+1$, and the final bound displayed at $p$ is
$F_p=\tfrac{\ell+1}4\bigl\{(1+\tfrac4\ell)\bar\dd-\tfrac\ell3\bar\rho+\tfrac4\ell\bigr\}$ \cite[pp.~28--29]{IUT4}, where
$\bar\dd$ and $\bar\rho$ are the averages over the two places. The values are in Table~\ref{tab:average}.
\end{theorem}

\begin{proof}
For a tuple $\bv$ and $t\in\{1,\dots,j\}$, \cite[Lemma~3]{Note} or (M$_{\mathrm g}$) at the tuple $\bv\circ(t\ j)$ and
\cite[Lemma~4]{Note} with the transposition $(t\ j)$, which fixes $0$, give an element of the component at $\bv$
whose factors have order at most $c(v_t)$ in total \cite[(1.1)--(1.3)]{Note}; so the hull component at $\bv$ has
log-volume at least $-\min_tc(v_t)$. For $B_j$: $\log(\dd^K_p)=\bar\dd\log p$, $\log(q_p)=2\ell\bar\rho\log p$ and
$\log(\mathfrak s^{\Q}_p)=\log p$ in the displays of \cite[pp.~28--29]{IUT4}. The values were computed exactly
(\S\ref{s:repro}).
\end{proof}

\begin{table}[ht]
\centering\small
\begin{tabular}{@{}lllll@{}}
\toprule
datum & $L_1-B_1$ & $L_2-B_2$ & $L_3-B_3$ & $\tfrac13\sum_jL_j-F_p$\\
\midrule
$D_2$ & $-\tfrac{727}{210}$ & $-\tfrac{517}{140}$ & $-\tfrac{1333}{420}$ & $-\tfrac{5477}{1470}$\\
$D_3$ & $-\tfrac{114}{35}$ & $-\tfrac{67}{70}$ & $\tfrac{1851}{280}$ & $\tfrac{3089}{5880}$\\
$D_{18}$ & $-\tfrac{153}{70}$ & $-\tfrac9{140}$ & $\tfrac{197}{35}$ & $\tfrac{2687}{2940}$\\
$D_{15}$ & $-\tfrac{16}7$ & $-\tfrac14$ & $\tfrac{31}7$ & $\tfrac{251}{588}$\\
\bottomrule
\end{tabular}
\caption{Averaged comparisons at $p$ (units of $\log p$). A positive entry means that the averaged lower bound for the
hull exceeds the displayed bound. For $D_2$ no averaged bound is exceeded; its non-containment is in a single
component.}\label{tab:average}
\end{table}

\begin{remark}
Two variants change the numbers but not the signs of the last column. If the members of the poly-isomorphism over
the transposition are taken to have identity components, as one may after identifying the members of the capsule with
the model of \cite[Def.~4.1(iii)]{IUT1}, the rounding $\lceil\cdot\rceil$ can be dropped; the last column
becomes $\tfrac{241}{245}$, $\tfrac{1059}{980}$, $\tfrac{249}{392}$ for $D_3$, $D_{18}$, $D_{15}$. If permutations
moving $0$ are admitted as well (L$^0$), the minimum runs over $t\in\{0,\dots,j\}$, and without rounding the last
column becomes $\tfrac{1217}{490}$, $\tfrac{4323}{1960}$, $\tfrac{1233}{784}$.
\end{remark}

\section{A global check for \texorpdfstring{$D_2$}{D2}}\label{s:global}

For $D_2$ the bad places are $b_1,b_2$ above $p$, the place above $R$ (the large prime factor of $\Nm(1-a)$, see
\cite[Prop.~2]{Note}) where $\operatorname{ord}(q)=2$, and the two places above $11$, where $\operatorname{ord}(q)=2$
each; so $\log q=\tfrac12(14\log p+2\log R+4\log11)=240.09\ldots$ in the normalization of \cite[Thm.~1.10]{IUT4}. The
final inequality of Theorem~1.10 holds for $D_2$ with a wide margin, its left side $\tfrac16\log q\approx40.0$ being
far below the constant $20e^*_{\mathrm{mod}}\ell=154828800$ on its right. The inequality that \cite{LANAiut}
formulates for Step~(xi), $(s-1)|\log q|\le\mulog(\mathrm{hull}(\mathrm{VC}(I)))$ with $s=\ell(\ell+1)/12$ and
$|\log q|=\log q/2\ell$ (file \nolinkurl{Iut/WallEquiv/README.md}), also holds for $D_2$: the left side is $62.88\ldots$,
while by the formula $\sum_c w_c\sum_j\rho_{c_j}\log p$ for the prime components recorded there, with
$\rho=1-1/e$ at the tame places above $p$ and $R$ and the other components nonnegative, the right side is at least
$3\bigl(\tfrac12(\tfrac{104}{105}+\tfrac{34}{35})\log p+\tfrac12\cdot\tfrac{104}{105}\log R\bigr)=221.0\ldots$.
The place above $R$, where $q$ has order only $2$, supplies most of this. Earlier, \cite{Takkun} (15~July 2026)
proved in Lean that, in its model, the hull of the possible images can be bounded by the theta volume plus a ledger
$L$ of the local terms of \cite[Props.~1.1--1.4]{IUT4}, compatibly with the containment of Step~(xi), if and only if
$(s-1)|\log q|\le L$ and the underlying model of Theorem~3.11 exists (\nolinkurl{thetaHullTranscription_iff}, file
\nolinkurl{IutLean/Multiradial/ThetaHullDichotomy.lean}; see also (F3) in
\nolinkurl{REQUEST_theta_hull_derivation.md}). This is the same inequality, with $L$ in place of
$\mulog(\mathrm{hull}(\mathrm{VC}(I)))$.

\section{Earlier work and dates}\label{s:prior}

\begin{itemize}
\item Explicit initial $\Theta$-data: \cite[\S2.10]{DH1} (11a1 over $\Q$, $\ell=13$); \cite{Zhou1,Zhou2} (families,
$\Fmod=\Q$); \cite[Def.~4.1]{MFHMP} ($\mu_6$-initial $\Theta$-data, no numerical example); initial $\Theta$-data
constructed in Lean from an elliptic curve and stated inputs, no numerical example: \cite{LANAiut}, file
\nolinkurl{Iut/Concrete/Existence.lean}, first committed on 2~September 2026 (UTC); an explicit example over an
extension of $\Q(i)$ was asked for in \cite{MO}.
\item The permutations of tensor factors in (Ind1): \cite[\S4]{DH1} (2020); \cite[\S6.2, footnote~7]{DH} (2020).
\item The two readings, the agreement when the orders of $q$ agree above each prime, the failure at two bad places with
different orders, the minimum versus the average, and primes with a good and a bad place from the Belyi reduction:
\cite{LANAiut}, file \nolinkurl{Iut/Concrete/Literal.lean}, first committed on 11~October 2026 (UTC); the per-image hull,
file \nolinkurl{Iut/Concrete/Ind1Invariance.lean}, the same day.
\item The question of the source of $\lambda=\operatorname{ord}(q^{j^2})$ in Step~(v), and the equivalence in Lean of the
hull bound with $(s-1)|\log q|\le L$ for a ledger $L$ (\S\ref{s:global}): \cite{Takkun} (15~July 2026). The last
estimate of Step~(v) restated without an independent proof: \cite[Prop.~6.10.9]{JoshiIV}.
\item A different loss in the hull of a tensor packet, from boxing the integral structure: \cite{Zeng} (first version 29~October 2026).
\item The data of this note: $D_{211}$ on 6 September 2026 \cite{First}, $D_{109}$ on 7 September 2026 \cite{Second},
$D_2$ on 6 October 2026 \cite{NoteV1}; $D_{2357}^{(\ell)}$ is the curve of \cite[\S13]{Second} (September 2026),
verified as an initial $\Theta$-datum for the fourteen values of $\ell$ in \cite{Fourteen} (6 November 2026); $D_3$,
$D_{18}$, $D_{15}$ and the values of \S\ref{s:comparisons} are new here.
\end{itemize}

\section{Reproducibility}\label{s:repro}

The values of Tables~\ref{tab:component} and~\ref{tab:average} are computed with exact rational arithmetic by two
standard-library Python programs: \texttt{catalogue.py} (Proposition~\ref{p:closed}, with an enumeration of all tuples
for $\ell=7$, $j\le3$, and checks against the values $\tfrac{27}{35}$, $\tfrac87$, $\tfrac{88}7$ recorded in
\cite{Note,First}) and \texttt{averaged.py} (Theorem~\ref{t:average}). The arithmetic inputs of
Proposition~\ref{p:new} are checked by the programs listed in Appendix~\ref{s:proofs}. All these programs, with
their output, are in the archive deposited with this note.

\appendix

\section{Proofs for the new data}\label{s:proofs}

In each case the proof of \cite[Prop.~2]{Note} and of its Table~1 applies with the changes below; in particular
\cite[Cor.~2.2]{IUT4} is not used, and conditions (d)--(f) are constructed directly from $j\notin\Q$, the split
multiplicative reduction at the selected places and the image containing $\mathrm{SL}_2(\F_7)$, as there.

\begin{proposition}\label{p:D3}
$D_3$: $a=\pi^6\bar\pi^{-30}$ with $\pi=14164+5825\sqrt5$.
\end{proposition}
\begin{proof}[Changes]
Now $\operatorname{ord}_{b_1}(j)=-12$ and $\operatorname{ord}_{b_2}(j)=-60$, so $j\notin\Q$. At $b_2$, where $a$ has a
pole of even order, the change of variables $x=aX$, $y=a\sqrt a\,Y$, with $\sqrt a=\pi^3\bar\pi^{-15}\in F_0$, gives
$Y^2=X(X-1)(X-a^{-1})$ with $\operatorname{ord}_{b_2}(a^{-1})=30$. The numerator of $1-a$ is
$-\prod_{d\mid6}\Phi_d(\pi,\bar\pi^{5})$, with the homogeneous cyclotomic forms $\Phi_d$, no two of which have a common
prime divisor outside $2,3$. None of the four norms $\Nm(\Phi_d(\pi,\bar\pi^5))$ ($38$, $38$, $75$ and $75$ digits) is
divisible by the seventh power of an odd prime (checked by a segmented sieve up to the seventh roots and independently
by a complete factorization with primality certificates), $1-a$ is a unit at the place above $3$ and has order $5$ at
the place above $2$. At the place above $7$, $a$ and $1-a$ are units. For the image, the place of degree one above
$41$ with $\sqrt5\equiv13$ gives $a_{41}=-6$, and $X^2+6X+41$ is irreducible over $\F_7$; the transvection comes from
$b_1$, where $\operatorname{ord}(q)=12$. At $b_1$ and $b_2$ the orders of $q$ are $12$ and $60$, so
$e=35$ and $7$, the differents have orders $34$ and $6$, and $\rho=\tfrac67$ and $\tfrac{30}7$.
\end{proof}

\begin{proposition}\label{p:D18}
$D_{18}$: $a=\pi^{18}$ with $\pi=14164+5825\sqrt5$.
\end{proposition}
\begin{proof}[Changes]
The six Legendre orbit comparisons of $a$ fail in $\Z[\sqrt5]$, so $j\notin\Q$. We have
$a-1=\prod_{d\mid18}\Phi_d(\pi)$; no norm of a factor is divisible by the seventh power of a prime (a sieve up to the
seventh roots of the norms, at most $2634850$), two factors meet only above $2$ and $3$, which are inert, the order of
$1-a$ at the place above $3$ is $0$, and the place above $2$ is not selected. So every positive order of $1-a$ at a
selected place is at most $6$, the only positive order of $a$ is $18$, at $v_\pi$, and $7$ divides neither
$\Nm(a)=p^{18}$ nor $\Nm(1-a)$. For the image, the place of degree one above $31$ with $\sqrt5\equiv6$ gives
$a_{31}=-8$, and $a_{31}^2-4\cdot31\equiv3\pmod7$ is not a square; the transvection comes from $v_\pi$, as $7\nmid36$.
Above $p$, $\bar\pi$ gives a good place $g$ with $e=1$; at $b=v_\pi$ the curve is split Tate over
$\Qp(\sqrt{-1})$ with $\operatorname{ord}_p(q)=36$, so $e=105/\gcd(105,36)=35$, the different has order $34$ and
$\rho=\tfrac{18}7$.
\end{proof}

\begin{proposition}\label{p:D15}
$D_{15}$: $a=\pi^{15}$ with $\pi=-5932-919\sqrt5$, $p=30965819$.
\end{proposition}
\begin{proof}[Changes]
$p$ is prime, $p\equiv4\pmod5$ and $p\equiv3\pmod4$; $p=\pi\bar\pi$ splits in $F_0$ and both places have degree one.
At $b=v_\pi$, $\operatorname{ord}(a)=15$; at $g=v_{\bar\pi}$, $a$ is a unit and $a\not\equiv1$; $1-a$ is a unit at both,
so $\operatorname{ord}_b(j)=-30$, $\operatorname{ord}_g(j)\ge0$ and $j\notin\Q$. The content of $a-1$ in
$\Z[(1+\sqrt5)/2]$ is $2$, and $5\nmid\Nm(a-1)$; hence at a split odd prime at most one place divides $a-1$, no
inert odd prime and not the ramified prime $5$ does, and at the place above $2$ (inert) $1-a$ has order $1$, where
$\operatorname{ord}(j)=6\ge0$. We have $\pi^{15}-1=\Phi_1(\pi)\Phi_3(\pi)\Phi_5(\pi)\Phi_{15}(\pi)$; the norms of the
blocks are pairwise coprime and none has a seventh-power prime factor other than at $2$ and $7$ (trial division up to
the seventh roots, at most $363942203$, by two methods with positive and negative controls), and $a-1$ has order $0$ at
the places above $3$ and $5$. So the orders of $1-a$ at selected places lie in $\{1,\dots,6\}$, the order of $a$ at $b$
is $15$, and $7$ divides neither $\Nm(a)$ nor $\Nm(1-a)$. For the image, at a place of degree one above $11$ the
reduction is good with $a_{11}=0$, and $X^2+4$ is irreducible over $\F_7$; the transvection comes from $b$, where
$\operatorname{ord}_w(q)=30e(w|b)$ is prime to $7$. At $b$ the curve is non-split multiplicative over $\Qp$ and split over
$\Qp(\sqrt{-1})$, with $\operatorname{ord}_p(q)=30$, so $e=7$, the different has order $6$ and $\rho=\tfrac{15}7$; at
$g$, $e=1$.
\end{proof}

The arithmetic inputs were checked by two programs for each datum, the second written independently from the
specification:
\begin{itemize}
\item $D_3$: \nolinkurl{certify_wide.py}, \nolinkurl{independent_check_wide.py};
\item $D_{18}$: \nolinkurl{certify_datum.py}, \nolinkurl{independent_check.py};
\item $D_{15}$: \nolinkurl{certify_mixed15.py}, \nolinkurl{independent_check_mixed15.py}.
\end{itemize}
For $D_3$, the complete factorization of the four norms, with primality certificates, is checked by
\nolinkurl{check_pratt_factorizations.py}.

\subsection*{Use of AI tools}
Generative AI tools (large language models: \mbox{Claude} by \mbox{Anthropic}, and \mbox{ChatGPT} and \mbox{Codex} by
\mbox{OpenAI}) were used extensively in this work, under the author's direction, in finding the data and the proofs,
in the computations and in the writing. The author takes full responsibility for the contents of this note.

\begin{thebibliography}{LANA-iut}

\bibitem[DH]{DH}
Taylor Dupuy and Anton Hilado, \emph{Probabilistic {S}zpiro, baby {S}zpiro, and explicit {S}zpiro from {M}ochizuki's
  {C}orollary 3.12}, 2020, arXiv:2004.13108v2 [math.NT].

\bibitem[DH1]{DH1}
Taylor Dupuy and Anton Hilado, \emph{The statement of {M}ochizuki's {C}orollary 3.12, initial theta data, and the
  first two indeterminacies}, 2020, arXiv:2004.13228 [math.NT]; \S4 and \S4.6 refer to version~1, \S2.10 to
  version~2 (2025).

\bibitem[First]{First}
Toshihisa Urahigashi, \emph{An explicit initial $\Theta$-datum with mixed reduction at a split prime}, version
  0.1.7, Zenodo, \href{https://doi.org/10.5281/zenodo.22537342}{doi:10.5281/zenodo.22537342}.

\bibitem[First-full]{FirstFull}
Toshihisa Urahigashi, \emph{An explicit initial $\Theta$-datum with mixed reduction over a split prime, and the local
  estimates of inter-universal {T}eichm\"uller theory~{IV}}, version 0.1.7, in the deposit of \cite{First}; its
  \S\S7--9 record a reading whose application to the source was withdrawn in that version.

\bibitem[Fourteen]{Fourteen}
Toshihisa Urahigashi, \emph{Fourteen initial $\Theta$-data on one curve over $\mathbb{Q}(i)$}, version 1.0.0, Zenodo,
  6 November 2026, \href{https://doi.org/10.5281/zenodo.23990717}{doi:10.5281/zenodo.23990717}.

\bibitem[Fourth]{Fourth}
Toshihisa Urahigashi, \emph{Split bad primes and the tuple structure of inter-universal {T}eichm\"uller theory:
  what the source states, and what follows (a fourth report)}, version 1.0.4, Zenodo,
  \href{https://doi.org/10.5281/zenodo.22648944}{doi:10.5281/zenodo.22648944}.

\bibitem[IUT1]{IUT1}
Shinichi Mochizuki, \emph{Inter-universal {T}eichm\"uller theory {I}: construction of {H}odge theaters}, Publ. Res.
  Inst. Math. Sci. \textbf{57} (2021), 3--207.

\bibitem[IUT3]{IUT3}
Shinichi Mochizuki, \emph{Inter-universal {T}eichm\"uller theory {III}: canonical splittings of the
  log-theta-lattice}, Publ. Res. Inst. Math. Sci. \textbf{57} (2021), 403--626.

\bibitem[IUT4]{IUT4}
Shinichi Mochizuki, \emph{Inter-universal {T}eichm\"uller theory {IV}: log-volume computations and set-theoretic
  foundations}, Publ. Res. Inst. Math. Sci. \textbf{57} (2021), 627--723; page numbers refer to the preprint
  version of April 2020.

\bibitem[JoshiIV]{JoshiIV}
Kirti Joshi, \emph{Construction of arithmetic {T}eichm\"uller spaces {IV}}, arXiv:2403.10430v2 [math.NT].

\bibitem[LANA-iut]{LANAiut}
lana-agents, \emph{iut}, Lean repository, \url{https://github.com/lana-agents/iut}, consulted at commit
  \texttt{ecb3c86} (1 November 2026).

\bibitem[MFHMP]{MFHMP}
Shinichi Mochizuki, Ivan Fesenko, Yuichiro Hoshi, Arata Minamide, and Wojciech Porowski, \emph{Explicit estimates in
  inter-universal {T}eichm\"uller theory}, Kodai Math. J. \textbf{45} (2022), no.~2, 175--236.

\bibitem[MO]{MO}
\emph{Explicit example of elliptic curve of the kind needed for {IUTT}}, MathOverflow question 299845 (2018).

\bibitem[Note]{Note}
Toshihisa Urahigashi, \emph{A local non-containment in Step~(v) of {IUT}~{IV}, Theorem~1.10}, version 3.0.0, Zenodo,
  1 November 2026, \href{https://doi.org/10.5281/zenodo.23852752}{doi:10.5281/zenodo.23852752}.

\bibitem[Note-v1]{NoteV1}
Toshihisa Urahigashi, \emph{A local non-containment in Step~(v) of {IUT}~{IV}, Theorem~1.10}, version 1.0.0, Zenodo,
  6 October 2026, \href{https://doi.org/10.5281/zenodo.23166433}{doi:10.5281/zenodo.23166433}.

\bibitem[Second]{Second}
Toshihisa Urahigashi, \emph{An initial $\Theta$-datum over an imaginary quadratic field, and numerical values of the
  local quantities of inter-universal {T}eichm\"uller theory~{IV}}, version 1.0.5, Zenodo,
  \href{https://doi.org/10.5281/zenodo.22601333}{doi:10.5281/zenodo.22601333}.

\bibitem[Takkun]{Takkun}
Takkun-kohinata, \emph{IUT\_LEAN}, GitHub repository, \url{https://github.com/Takkun-kohinata/IUT_LEAN}, files
  \nolinkurl{REQUEST_theta_hull_derivation.md} and \nolinkurl{IutLean/Multiradial/ThetaHullDichotomy.lean}; commit
  \texttt{9d56c46} (15 July 2026, UTC).

\bibitem[Zeng]{Zeng}
Chunping Zeng, \emph{An integral-structure inflation in the hull operation of {IUT} {C}orollary 3.12}, Zenodo,
  version of 30 October 2026, \href{https://doi.org/10.5281/zenodo.23811410}{doi:10.5281/zenodo.23811410}; first
  version 29 October 2026; all versions \href{https://doi.org/10.5281/zenodo.23754276}{doi:10.5281/zenodo.23754276}.

\bibitem[Zhou1]{Zhou1}
Zhong-Peng Zhou, \emph{The inter-universal {T}eichm\"uller theory and new {D}iophantine results over the rational
  numbers.~{I}}, arXiv:2503.14510 [math.NT], 2025.

\bibitem[Zhou2]{Zhou2}
Zhong-Peng Zhou, \emph{The inter-universal {T}eichm\"uller theory and new {D}iophantine results over rational
  numbers.~{II}}, arXiv:2510.05448 [math.NT], 2025.

\end{thebibliography}

\end{document}
